\chapter{CHALLENGES AND SOLUTIONS}

\section{Challenge 1}

\textbf{Problem.}
The first challenge was defining a project scope that was realistic for the
internship period while still covering the browser extension, study web app,
and AWS backend.

\textbf{Resolution.}
The scope was narrowed to an offline-first flashcard workflow with local
capture, cloud sync, study, export, and cleanup.

\textbf{Lesson.}
Clear scope definition early in the project prevents later rework and keeps the
system design coherent.

\section{Challenge 2}

\textbf{Problem.}
The second challenge was getting the local extension flow and the remote AWS
API flow to work together reliably.

\textbf{Resolution.}
The project used a local storage fallback, JWT-protected sync endpoints, and
step-by-step verification of the capture, sync, and study paths.

\textbf{Lesson.}
End-to-end validation is necessary when a feature spans the browser, backend,
and cloud infrastructure.

\section{Challenge 3}

\textbf{Problem.}
The third challenge was keeping the documentation aligned with the evolving
implementation and project closeout timeline.

\textbf{Resolution.}
The worklog, workshop pages, proposal, and final report were updated together
so the written material matched the verified project state.

\textbf{Lesson.}
Documentation should be treated as part of the system, not as an afterthought.

\section{Differences Between Documentation and Running Work}

Where the written plan and the running implementation differed, I used the
current worklog and workshop content as the source of truth instead of older
drafts.

\section{Confidentiality of Evidence}

The evidence included in the report excludes credentials, tokens, personal
data beyond the student information already provided, customer data, and other
confidential company information.